\documentclass[tikz,border=3mm]{standalone}
% Spolecny styl vlastnich vektorovych ilustraci; preklad jednotlivych stackup_*.tex.
\usepackage[czech]{babel}
\usepackage[utf8]{inputenc}
\usepackage[T1]{fontenc}
\usepackage{tgheros}
\renewcommand{\familydefault}{\sfdefault}
\definecolor{sig}{RGB}{244,185,75}
\definecolor{gnd}{RGB}{107,183,141}
\definecolor{pwr}{RGB}{228,129,112}
\definecolor{frfour}{RGB}{245,237,214}
\tikzset{copper/.style={draw=black!55,line width=0.4pt},
 layerlabel/.style={anchor=east,font=\small},
 note/.style={font=\small,align=center}}

\begin{document}
\begin{tikzpicture}
 \fill[frfour] (0,0.25) rectangle (8,3.25);
 \draw[copper,fill=sig] (0,3.25) rectangle (8,3.6);
 \node at (4,3.425) {\textbf{SIG}: signály -- reference L2 (GND)};
 \node[layerlabel] at (-0.15,3.425) {L1 / F.Cu};
 \node[font=\scriptsize] at (4,3.025) {Tenké dielektrikum (prepreg)};
 \draw[copper,fill=gnd] (0,2.45) rectangle (8,2.8);
 \node at (4,2.625) {\textbf{GND}: souvislá referenční plocha};
 \node[layerlabel] at (-0.15,2.625) {L2 / In1.Cu};
 \node[note] at (4,1.75) {Obvykle silnější jádro FR4\\Kapacita GND--PWR nemusí být velká.};
 \draw[copper,fill=pwr] (0,0.7) rectangle (8,1.05);
 \node at (4,0.875) {\textbf{PWR}: napájecí plocha nebo více oblastí};
 \node[layerlabel] at (-0.15,0.875) {L3 / In2.Cu};
 \node[font=\scriptsize] at (4,0.475) {Tenké dielektrikum (prepreg)};
 \draw[copper,fill=sig] (0,-0.1) rectangle (8,0.25);
 \node at (4,0.075) {\textbf{SIG} -- nejbližší reference L3 (PWR)};
 \node[layerlabel] at (-0.15,0.075) {L4 / B.Cu};
 \node[font=\scriptsize,anchor=west] at (0,-0.55) {Schématický řez běžné geometrie; konkrétní tloušťky určuje výrobce.};
\end{tikzpicture}
\end{document}
